% !TeX root = simplesnt-doc.tex
\documentclass{simplesnt}

%% Helpful packages
\usepackage{simplebnf}[2023/11/25]
\usepackage{kotex}
\usepackage{xcolor}
\usepackage{amsmath,amssymb}
\usepackage{amsfonts}
\usepackage{tikz}
\usetikzlibrary{calc}
\usepackage{siunitx}
\usepackage{caption}
\usepackage{listings}
\usepackage{mathrsfs}
\usepackage{stmaryrd}
\usepackage{subcaption}
\usepackage{graphicx}
\usepackage{colortbl}


\captionsetup{font=small}
\lstset{basicstyle = \small\ttfamily, breaklines}

\usepackage{tabularray}
\UseTblrLibrary{booktabs}

\usepackage{metalogo}
\usepackage{kotex-logo}


\newcommand{\textalt}[2]{#1\textsubscript{\ensuremath{\text{\scriptsize\textit{#2}}}}}
\newcommand{\CILmm}{\texorpdfstring{CIL-\kern0.08em-}{CIL--}}
\newcommand{\fainttablerow}{\\[0.2ex]\arrayrulecolor{black!12}\hline\arrayrulecolor{black}}
\newcommand*{\presentfont}[1]{{\fontspec{#1}#1}}
\ExplSyntaxOn
\bool_if:nF { \sys_if_engine_xetex_p: || \sys_if_engine_luatex_p: }
  { \newcommand*{\fontspec}[1]{} }
\ExplSyntaxOff

\title{\CILmm{} 합성을 통한 Sparrow 분석기 공격}
\author{정원준}

\begin{document}
\date{2026년 7월 10일}
\maketitle


\begin{frame}
  \frametitle{목차}
  \tableofcontents
\end{frame}



\section{Sparrow 분석기 공격 목표}


\begin{frame}
  \frametitle{Sparrow 공격 목표}
  \begin{itemize}
    \uncover<1->{
      \item completeness 공격: 
        \newline abstract semantics가 concrete semantics를 너무 크게 포함함
        \newline (concrete = abstract인 그림) (concrete < abstract인 그림) (abstract = top인 그림)
    }
    \vspace{2em}
    \uncover<2->{
      \item soundness 공격: 
        \newline concrete semantics가 abstract semantics에 포함되지 않음(버그)
        \newline (concrete $\notin$ abstract인 그림)
    }
  \end{itemize}
\end{frame}


\begin{frame}
  \frametitle{Sparrow의 기본 구조}
  (C->CIL->CFG 의 예시 그림)

  예시는 

  int main() \{ int x = 1; if(x == 1) { x++; } return x; \}

  정도로
\end{frame}


\section{\CILmm{} 언어}

\begin{frame}
  \frametitle{CIL 언어}
  CIL이 뭐냐, C의 부분집합 느낌. 분석을 쉽게 할 수 있는 언어
  애바하거나 복잡한 구조 단순화, side effect 제거, 정규화
\end{frame}

\begin{frame}
  \frametitle{CIL-- 언어}
  completeness를 공격하기 위해서는 튜링 완전한 subset이기만 하면 충분
  soundness를 공격하기 위해서는 많은 feature가 있을수록 좋음(버그 찾는거니까)

  CIL- : 의미적으로 별로 필요없는것들 제거 (typedef, struct 등)
  CIL-- : 나중에 추가했으면 좋겠지만 일단 미니멀하게 가기 위해 제거 (포인터 + malloc/free, 배열, int 외 type)
\end{frame}

\section{실행의미를 Big-Step으로 정의}

\begin{frame}
  \frametitle{으에}
  CompCert의 Clight Big-Step 의미구조가 있긴한데, 여기서는 CIL--에서의 의미구조를 직접 정의

  expression은? 메모리를 보면서 정의
  메모리는? stack(frame)과 heap이 있는데, 아직 heap은 없음

  call은? frame을 하나 넣고, arg랑 로컬 변수 할당, return하는거까지 표현 (그림 넣기)

\end{frame}

\begin{frame}
  \frametitle{으에}
  동치성
  ast 동치, bigstep tree 동치, 메모리 동치
  이게 언제 필요하지?
\end{frame}

\section{Sparrow의 실행의미와 대응시키기}

\begin{frame}
  \frametitle{으에}
  C->CIL->CFG
     ->CIL-- 구조.
  C쪽은 더 이상 볼 필요가 없음

  가장 큰 문제: sparrow는 concrete semantics를 정확히 정의하지 않음. 

  C: ISO C 표준이 있는데, formal하진 않음
  CIL: 그런것도 없음. 
  Sparrow: 얘도 아예 없음
  공격기: CIL에 대한 Big-Step tree랑, derivator가 있긴 함

  공격기의 concrete 실행의미가 Sparrow의 abstract 실행의미에 포함되지 않더라도, soundness가 깨졌다! 라고 말할 수 있는가?
  (여기에 그림 넣기)

  
  C 언어의 비결정성 보기

  각종 비결정성을 모두 제거해야 공격을 제대로 할 수 있음
\end{frame}


\begin{frame}
  \frametitle{C 언어의 비결정성}

  {\scriptsize
  \centering
  \setlength{\tabcolsep}{2.0pt}
  \renewcommand{\arraystretch}{1.08}
  \begin{tabular}{@{}p{0.22\linewidth}p{0.31\linewidth}p{0.28\linewidth}p{0.16\linewidth}@{}}
    \toprule
    \textbf{예} & \textbf{C} & \textbf{CIL} & \textbf{\CILmm{}} \\
    \midrule
    \texttt{h(f(),g())} \newline \texttt{f()+g()} \newline \texttt{\{f(),g()\}} & 계산 순서 비결정 & 함수 호출은 expression으로 들어가지 않음 & OK \fainttablerow
    \texttt{sizeof(int[n++])} & \texttt{n++}을 실행하거나/하지 않거나 & \texttt{n++} 문법이 없음 & OK \fainttablerow
    \texttt{i = i++ + 1;} & undefined behavior & \texttt{i++} 문법이 없음 & OK \fainttablerow
    \texttt{int x; use(x);} & indeterminate value & indeterminate value & 합성 시 생성 x\fainttablerow
    \texttt{"a" == "a"} & true/false 모두 가능 & 여전히 비결정적 & string 없음 \fainttablerow
    \texttt{char c = -1;} & -1/255 & 여전히 비결정적 & char 없음 \fainttablerow
    \texttt{x >> 1} & \texttt{x} $< 0$ 일 때, & 여전히 비결정적 & 비트연산 없음 \fainttablerow
    \texttt{malloc(n)} & 메모리 주소가 비결정적 & fresh object & fresh object \\
    \bottomrule
  \end{tabular}
  }

  \uncover<2->{
    이외에도 malloc 실패 여부, struct의 byte 패딩, 정수 overflow 등 여러 가지 고려해야 할 비결정성이 많음
  }
\end{frame}

\section{메모리를 포함하여 합성하는 방법}

\begin{frame}
  \frametitle{으에}
  그러게 어떻게 하지...
  증명나무 조각들을 레고처럼 끼워서 이어붙이려면 sub-prooftree의 모든 구조가 다 제대로 있어야 함.
  즉, 전/후 메모리가 잘 정의되어 있어야 함
  같은 expression "x"여도, 메모리 하나당 한 개의 proof tree가 생기고 이들은 서로 구분됨.
\end{frame}

\section{앞으로 할 일}

\begin{frame}
  \frametitle{으에}
  syntax, semantics(big-step), interpreter 등등 기본적인거 짜 놨으니
  이제 본격적으로 합성 시작.

  걱정거리: simple.c만 해도 크기가 매우매우 큼. (그림 삽입)
  재귀적으로 합성할 것/아닌 것을 구분해서 쉬운 step은 바로바로 진행되게끔 하는 스킬이 필요할 듯
\end{frame}



% 감사합니다

\definecolor{mydarkblue}{RGB}{32,56,61}

\setbeamercolor{background canvas}{bg=mydarkblue}
\setbeamercolor{normal text}{fg=white}
\setbeamercolor{frametitle}{fg=white}
\addtocounter{framenumber}{-1} % 마지막 슬라이드 번호를 하나 줄임
\appendix
\begin{frame}[c] % [c]는 내용 수직 가운데 정렬
  \centering
  {\Large \textcolor{white}{감사합니다}}
\end{frame}

\end{document}
